% Pista geo-26s-q4b · Geometria
% Procedència: PAU Catalunya 2026 · Convocatòria de setembre · Sèrie 2 · Exercici 4 — Opció B
\documentclass[11pt,a4paper]{article}
\usepackage{pista}
\pistainfo{Catalunya 2026}{Convocatòria de setembre}{Sèrie 2}

\begin{document}

\section*{\textcolor{blauPista}{Exercici 4 --- Opció B}}

\noindent L'Aina vol construir una capsa de cartró amb forma de prisma rectangular amb tapa, per guardar-hi la tauleta i el carregador.

\noindent L'Aina vol fer un croquis en 3D de la capsa utilitzant un programa informàtic. Per a començar, situa un costat de la tapa des del vèrtex $A(1, 1, 1)$ fins al vèrtex $B(4, 1, 5)$, i l'altre costat de la tapa des del vèrtex $A$ fins al vèrtex $C(1, 10, 1)$.

\subsection*{4.1. \normalfont Quines són les coordenades del quart vèrtex de la tapa? Calculeu les longituds dels dos costats de la tapa. \hfill\textnormal{[1~punt]}}

\pista{Fixa't que els dos costats surten d'$A$: $B$ i $C$ són els veïns d'$A$, i el quart vèrtex $D$ és l'oposat a $A$. Per arribar a $D$, surt de $B$ i aplica-hi el mateix desplaçament que porta d'$A$ a $C$: $D = B + \vec{AC}$. L'error habitual és tractar $BC$ com un costat (per exemple, fent $A + \vec{BC}$), quan en realitat és una diagonal. Si vols, comprova que $\vec{AB} \cdot \vec{AC} = 0$ (angle recte). Les longituds dels costats són els mòduls $|\vec{AB}|$ i $|\vec{AC}|$.}

\subsection*{4.2. \normalfont Sabent que els altres quatre vèrtexs de la capsa estan situats al pla $\pi: -4x + 3z = 24$, calculeu les coordenades d'un d'aquests vèrtexs (qualsevol) i l'alçària de la capsa. \hfill\textnormal{[1,5~punts]}}

\pista{La capsa és un prisma recte: la base és paral·lela a la tapa i cada vèrtex de la base és la projecció ortogonal d'un vèrtex de la tapa sobre $\pi$. Si vols, comprova que $\pi$ és paral·lel a la tapa ($\vec{AB} \times \vec{AC}$ ha de ser proporcional a $\vec{n}_\pi = (-4, 0, 3)$). Aleshores l'alçària és la distància d'un vèrtex de la tapa, per exemple $A$, al pla $\pi$ (escriu-lo abans com $-4x + 3z - 24 = 0$). Per al vèrtex, projecta $A$ ortogonalment sobre $\pi$: recta per $A$ amb vector director $\vec{n}_\pi$, i interseca-la amb $\pi$. Compte: un punt qualsevol de $\pi$ no serveix.}

\end{document}
